\begin{theorem}[The quadratic sector along $\Delta$]
\label{thm:c9-quadratic}
In the sectors $(A|_{W_{\mathbb R}^\perp},A^*|_{W_{\mathbb R}^\perp})$, $(A|_V,A^*|_V)$, $(\lambda,\lambda^*)$ and $(c,c^*)$ of degree $-1$ and the corresponding sectors $(A,A)$, $(\lambda,\lambda)$, $(c,A^*|_V)$ of degree $0$, with $x$ the position, write the sector parts of $S^{(1)}(\Delta)$ as
\[
S^{(1)}(\Delta)_\perp=\int\big(2E(A|_{W_{\mathbb R}^\perp})+2A|_{W_{\mathbb R}^\perp},A^*|_{W_{\mathbb R}^\perp}\big),\qquad S^{(1)}(\Delta)_\lambda=-\int\big(2E(\lambda)+3\lambda,\lambda^*\big),\qquad S^{(1)}(\Delta)_V=\int\big(2E(A|_V)+2A|_V,A^*|_V\big)-\int\big(2E(c),c^*\big),
\]
whose sum is the quadratic part of $S^{(1)}(\Delta)$; each is $Q$-closed.
\begin{itemize}
\item (i) The quadratic part of the arity-one counterterm is
\[
I^{CT}_{\mathfrak{psl}(4|4)}[\varepsilon]^{(1)}(\Delta)=\frac{\hbar\operatorname{tr}_{\mathfrak g}(\mathrm{ad}_X\mathrm{ad}_Y)}{16\pi^2(X,Y)}\log\varepsilon\Big[\int(d(\iota_xA|_V),A^*|_V)-\int(\iota_xdc,c^*)-\int(\rho(x){\not d}\lambda,\lambda^*)\Big]=\frac{\hbar\operatorname{tr}_{\mathfrak g}(\mathrm{ad}_X\mathrm{ad}_Y)}{16\pi^2(X,Y)}\log\varepsilon\;Q\Big(\int(\iota_xA|_V,c^*)-\tfrac12\int(\rho(x)\lambda^*,\lambda^*)\Big).
\]
Its coefficient of $\varepsilon^{-1}$ vanishes, and so does its part in the sector $(A|_{W_{\mathbb R}^\perp},A^*|_{W_{\mathbb R}^\perp})$.
\item (ii) The quadratic part of the arity-one obstruction, the sum \eqref{eq:c9-local} over the defect cycles with two legs, is
\begin{equation}
\hbar\,\mathrm{Ob}^{(1)}(\Delta)=\frac{\hbar\operatorname{tr}_{\mathfrak g}(\mathrm{ad}_X\mathrm{ad}_Y)}{16\pi^2(X,Y)}\Big[-\tfrac{11}3\int(dA,dA)+\tfrac{47}3\int(\lambda,{\not d}\lambda)-\tfrac{14}3\int(dc,A^*)\Big],
\label{eq:c9-ObDelta}
\end{equation}
$\int(dA,dA)=\int(dA|_V,dA|_V)+\int(dA|_{W_{\mathbb R}^\perp},dA|_{W_{\mathbb R}^\perp})$, with vanishing coefficient of $\int(d^*A|_V,d^*A|_V)$ and of every functional with an explicit factor $x$. The defect cycles with $S^{(1)}(\Delta)$ at the special edge give $\big[-4\int(dA,dA)+16\int(\lambda,{\not d}\lambda)-4\int(dc,A^*)\big]$ units and those with $S^{(1)}(\Delta)$ on a line give $\big[\tfrac13\int(dA,dA)-\tfrac13\int(\lambda,{\not d}\lambda)-\tfrac23\int(dc,A^*)\big]$ units: placement class by placement class these are the first summand and the composite part of \eqref{eq:c9-Ob}.
\item (iii) The quadratic part of \eqref{eq:c9-arity1} is solved exactly by
\begin{equation}
\begin{aligned}
\hbar F_1^{(1)}(\Delta)=\ &\frac{\hbar\operatorname{tr}_{\mathfrak g}(\mathrm{ad}_X\mathrm{ad}_Y)}{16\pi^2(X,Y)}\Big[\tfrac{22}3\int(A,A^*)-\tfrac{47}3\int(\lambda^*,\lambda)-\tfrac83\int(c^*,c)\Big]+\text{units}\cdot\Big(\theta\,S^{(1)}(\Delta)_\perp+\theta_\lambda\,S^{(1)}(\Delta)_\lambda+\theta_V\,S^{(1)}(\Delta)_V\Big)\\
&+\text{units}\cdot Q\Big(\kappa_V\int(\iota_xA|_V,c^*)+\kappa_\lambda\int(\rho(x)\lambda^*,\lambda^*)\Big),
\end{aligned}
\label{eq:c9-F1Delta}
\end{equation}
with free numbers $\theta,\theta_\lambda,\theta_V,\kappa_V,\kappa_\lambda$. The first bracket is the unique translation-invariant solution; it is $4$ units times $2\int(A,A^*)-4\int(\lambda^*,\lambda)-\int(c^*,c)$ plus $\big[-\tfrac23\int(A,A^*)+\tfrac13\int(\lambda^*,\lambda)+\tfrac43\int(c^*,c)\big]$ units. On the reflection-even functionals the free directions are all the $Q$-closed quadratic local functionals of degree $-1$, Weyl weight $4$ and degree at most one in $x$ (Proposition \ref{prop:c11-tc0}); modulo $Q$-exact functionals, the class of $\hbar F_1^{(1)}(\Delta)$ minus the first bracket is $\theta[S^{(1)}(\Delta)_\perp]+\theta_\lambda[S^{(1)}(\Delta)_\lambda]+\theta_V[S^{(1)}(\Delta)_V]$ units. In the sector $(A|_{W_{\mathbb R}^\perp},A^*|_{W_{\mathbb R}^\perp})$ the solutions are
\begin{equation}
\hbar F_1^{(1)}(\Delta)=\frac{\hbar\operatorname{tr}_{\mathfrak g}(\mathrm{ad}_X\mathrm{ad}_Y)}{16\pi^2(X,Y)}\Big[\tfrac{22}3\int(A|_{W_{\mathbb R}^\perp},A^*|_{W_{\mathbb R}^\perp})+\theta\int\big(2E(A|_{W_{\mathbb R}^\perp})+2A|_{W_{\mathbb R}^\perp},A^*|_{W_{\mathbb R}^\perp}\big)\Big],\qquad\theta\ \text{a number}.
\label{eq:c9-Fperp}
\end{equation}
\end{itemize}
The numbers $\theta,\theta_\lambda,\theta_V$ are fixed by the arity-two equations to $-\tfrac13$ (Theorem \ref{thm:c9-theta2}).
\end{theorem}

\begin{proof}
(i) By Proposition \ref{prop:c9-structure} (i), the counterterm is the singular part of the composite cycles with two legs, in which the composite kernel closes two legs of a tree of $I$ with four legs; their exact evaluation gives the stated multiple of $\log\varepsilon$, with vanishing coefficient of $\varepsilon^{-1}$, and the graphs with $S^{(1)}(\Delta)$ joined by an edge to the one-loop two-point subgraphs and to the vertex $-I^{CT}[\varepsilon]$ have opposite singular parts. In the sector $(A|_{W_{\mathbb R}^\perp},A^*|_{W_{\mathbb R}^\perp})$ the exact evaluation of the composite cycle gives the coefficient $0$ of $\log\varepsilon$. The second form follows from the substitutions of Proposition \ref{prop:c1-Q}: $Q\int(\iota_xA|_V,c^*)=-\int(\iota_xdc,c^*)+\int(d(\iota_xA|_V),A^*|_V)$, and $-\tfrac12Q\int(\rho(x)\lambda^*,\lambda^*)=-\int(\rho(x){\not d}\lambda,\lambda^*)$ by the symmetry of $(\rho(x)\zeta_1,\zeta_2)$. (ii) By Proposition \ref{prop:c9-structure} (ii) with $n=2$, whose hypothesis is Theorem \ref{thm:c4-finite}: the defect cycles with two legs have one or two vertices of $I$, the vertex $S^{(1)}(\Delta)$ at the special edge or on a line, and are evaluated exactly by their simplex forms, with the position realised by a momentum injected at $S^{(1)}(\Delta)$; the cycles with $S^{(1)}(\Delta)$ at the special edge and those with $S^{(1)}(\Delta)$ on a line give the two stated brackets, and the coefficients of the functionals with an explicit factor $x$ cancel; by Proposition \ref{prop:c11-position} (i) this is the vanishing of the quadratic part of $\mathrm{Ob}^{(1)}(P_b)$, which the defect cycles along $P_b$ give directly (Theorem \ref{thm:c9-P1} (i)). By Proposition \ref{prop:c9-structure} (iii) the same functional is $4$ times the quadratic part of Theorem \ref{thm:c4-quadratic} plus the image under $Q$ of the renormalised composite two-point functional of Remark \ref{rem:c9-Pi}; the two evaluations agree sector by sector. (iii) By Proposition \ref{prop:c1-Q}, $Q\int(A|_{W_{\mathbb R}^\perp},A^*|_{W_{\mathbb R}^\perp})=-\tfrac12\int(dA|_{W_{\mathbb R}^\perp},dA|_{W_{\mathbb R}^\perp})$, $Q\int(A|_V,A^*|_V)=-\tfrac12\int(dA|_V,dA|_V)-\int(dc,A^*)$, $Q\int(\lambda^*,\lambda)=-\int(\lambda,{\not d}\lambda)$ and $Q\int(c^*,c)=-\int(dc,A^*)$. These four functionals span the parity-even translation-invariant local functionals of degree $-1$ and Weyl weight $4$ in the four sectors, and their images are linearly independent, which gives existence and uniqueness of the translation-invariant solution; $Q$ of the first bracket is \eqref{eq:c9-ObDelta}. The sector parts are $Q$-closed because $v.$ commutes with $Q$ and $Q$ preserves the sectors $(A|_{W_{\mathbb R}^\perp},A^*|_{W_{\mathbb R}^\perp})$ and $(\lambda,\lambda^*)$ and mixes $(A|_V,A^*|_V)$ with $(c,c^*)$ only; Proposition \ref{prop:c11-tc0} gives the closed and the exact directions. In the sector $(A|_{W_{\mathbb R}^\perp},A^*|_{W_{\mathbb R}^\perp})$ the functionals of degree at most one in $x$ are spanned by $\int(A|_{W_{\mathbb R}^\perp},A^*|_{W_{\mathbb R}^\perp})$ and $\int(E(A|_{W_{\mathbb R}^\perp}),A^*|_{W_{\mathbb R}^\perp})$, and $Q\int(E(A|_{W_{\mathbb R}^\perp}),A^*|_{W_{\mathbb R}^\perp})=\tfrac12\int(dA|_{W_{\mathbb R}^\perp},dA|_{W_{\mathbb R}^\perp})$; the equation fixes the difference of the two coefficients to $\tfrac{22}3$.
\end{proof}
